\section*{5 Bit pigeonhole principle with parity gadget}
In this section, we prove the following theorem.
\begin{authorclaim}{Theorem 3}
Let $\ell$ and $k$ be natural numbers such that $2\le k\le\ell-7$. Then
\[
R^{1/3}_{\mathrm{pub}}(\oplus_k\operatorname{Search}(\mathrm{BPHP}_{2^\ell}^{2^\ell+2^k}))
=\Omega\left(\frac{2^{\ell/2}}{k2^{3k/2}}\right).
\]
For $k=2$ the stronger bound holds:
$R^{1/3}_{\mathrm{pub}}(\oplus_2\operatorname{Search}(\mathrm{BPHP}_{2^\ell}^{2^\ell+4}))=\Omega(2^\ell)$.
\end{authorclaim}
We consider a combinatorial analogue of the communication problem $\oplus_k\operatorname{Search}(\mathrm{BPHP}_{2^\ell}^m)$. Assume that each of $k$ parties gets $m$ binary strings from $\{0,1\}^\ell$, where $m>2^\ell$. The $i$th party has numbers $a_{i,1},\ldots,a_{i,m}\in\{0,1\}^\ell$ on their forehead. Based on these strings we form $m$ vectors from $\mathbb F_2^\ell$: $x_1,\ldots,x_m$, where $x_j=\sum_{i=1}^ka_{i,j}$. The goal is to find different $t,s\in[m]$ such that $x_t=x_s$. We denote this problem by $\oplus_k\mathrm{BPHP}_{2^\ell}^m$. It is straightforward that $R^{1/3}_{\mathrm{pub}}(\oplus_k\operatorname{Search}(\mathrm{BPHP}_{2^\ell}^m))\ge R^{1/3}_{\mathrm{pub}}(\oplus_k\mathrm{BPHP}_{2^\ell}^m)$, hence it is sufficient to prove a lower bound on the latter.
\begin{authorclaim}{Theorem 18}
Let $\ell,k$ be natural numbers such that $2\le k\le\ell-7$. Then
\[
R^{1/3}_{\mathrm{pub}}(\oplus_k\mathrm{BPHP}_{2^\ell}^{2^\ell+2^k})
=\Omega\big(R^{1/3}_{\mathrm{pub}}(\mathrm{UDISJ}_{k,2^{\ell-k}-1})-\ell\big).
\]
\end{authorclaim}
\begin{authorclaim}{Corollary 19}
$R^{1/3}_{\mathrm{pub}}(\oplus_k\mathrm{BPHP}_{2^\ell}^{2^\ell+2^k})=\Omega(2^{\ell/2}/(k2^{3k/2}))$. For $k=2$: $R^{1/3}_{\mathrm{pub}}(\oplus_2\mathrm{BPHP}_{2^\ell}^{2^\ell+4})=\Omega(2^\ell)$.
\end{authorclaim}
\begin{proof}[Proof of Corollary 19]
Follows from Theorem 18 and Theorem 7; for $k=2$ apply Theorem 8.
\end{proof}
Theorem 3 immediately follows from Corollary 19.
\subsection*{5.1 Warm-up example}
Consider $\mathrm{Distinct}_{k,\ell}$: each of $k$ parties has a matrix from $\mathbb F_2^{2^\ell\times\ell}$ on their forehead. The goal is to determine whether all rows of the sum of these matrices are distinct. A version without the xor-gadget is referred to as Element Distinctness (ED) [25].
\begin{authorclaim}{Proposition 20}
$R^{1/3}_{\mathrm{pub}}(\mathrm{Distinct}_{k,\ell})\ge R^{1/3}_{\mathrm{pub}}(\mathrm{UDISJ}_{k,2^{\ell-k}})$.
\end{authorclaim}
Let $S_k$ denote the set of matrices from $\{0,1\}^{2^k\times k}$ with all distinct rows. Let $K_k\in\{0,1\}^{2^k\times k}$ have $i$th row $\operatorname{bin}_k(i-1-((i-1)\bmod2))$, i.e. rows $\operatorname{bin}_k(0),\operatorname{bin}_k(0),\operatorname{bin}_k(2),\operatorname{bin}_k(2),\ldots,\operatorname{bin}_k(2^{k-1}-2),\operatorname{bin}_k(2^{k-1}-2)$. Every row starts with zero and appears exactly twice.
The following combinatorial lemma is proved in Subsection 5.4.
\begin{authorclaim}{Lemma 21}
There exist matrices $A_1(0),A_1(1),\ldots,A_k(0),A_k(1)\in\mathbb F_2^{2^k\times k}$ such that $\sum_{i=1}^kA_i(1)=K_k$ and for all $b_1,\ldots,b_k\in\{0,1\}$, if $\bigwedge_{i=1}^kb_i=0$, then $\sum_{i=1}^kA_i(b_i)\in S_k$.
\end{authorclaim}
\begin{proof}[Proof of Proposition 20]
Let $(x_{i,1},\ldots,x_{i,2^{\ell-k}})$ be an input of the $i$th party of the problem $\mathrm{UDISJ}_{k,2^{\ell-k}}$. For all $i\in[k]$ we construct a matrix $D_i$ of size $2^\ell\times\ell$ and put it on the forehead of the $i$th party. Let $A_i(b)$ for $i\in[k]$, $b\in\{0,1\}$ be matrices of size $2^k\times k$ from Lemma21. Let $J_t$ for $t\in[1,\ldots,2^{\ell-k}]$ be a matrix of size $2^k\times(\ell-k)$ such that all its rows are equal to $\operatorname{bin}_{\ell-k}(t-1)$. Let us define
\[
D_1:=\begin{pmatrix}J_1&A_1(x_{1,1})\\\vdots&\vdots\\J_j&A_1(x_{1,j})\\\vdots&\vdots\\J_{2^{\ell-k}}&A_1(x_{1,2^{\ell-k}})\end{pmatrix};
\quad
D_i:=\begin{pmatrix}0_{2^k\times(\ell-k)}&A_i(x_{i,1})\\\vdots&\vdots\\0_{2^k\times(\ell-k)}&A_i(x_{i,j})\\\vdots&\vdots\\0_{2^k\times(\ell-k)}&A_i(x_{i,2^{\ell-k}})\end{pmatrix}
\quad(i\in\{2,\ldots,k\}).
\]
By Lemma21, the matrix $D_1+D_2+\cdots+D_k$ has the following property: for all $j\in[2^{\ell-k}]$, its submatrix formed by the rows with numbers from $[2^k\cdot(j-1)+1,2^k\cdot j]$ has two equal rows if and only if $x_{1,j}=x_{2,j}=\cdots=x_{k,j}=1$. Thus, the communication complexity of $\mathrm{UDISJ}_{k,2^{\ell-k}}$ is at most the communication complexity of $\mathrm{Distinct}_{k,\ell}$.
\end{proof}
\subsection*{5.2 Proof of Theorem 18}
We modify Proposition20 to reduce $\mathrm{UDISJ}_{k,2^{\ell-k}-1}$ to $\oplus_k\mathrm{BPHP}_{2^\ell}^{2^\ell+2^k}$ by adding fake rows not corresponding to the disjointness input and randomizing to hide them among other rows.
\begin{proof}[Proof of Theorem 18]
For $N>2^\ell$, define $\mathrm{ROW}\oplus_k\mathrm{BPHP}_{2^\ell}^N$: party $i$ has $M_i\in\mathbb F_2^{N\times\ell}$ on their forehead and the goal is to find the value of a row appearing at least twice in $M_1+\cdots+M_k$, rather than the indices of equal rows.
\begin{authorclaim}{Claim 22}
If $R^{1/3}(\oplus_k\mathrm{BPHP}_{2^\ell}^N)\le t$, there exists a protocol $P$ for $\mathrm{ROW}\oplus_k\mathrm{BPHP}_{2^\ell}^N$ using $O(t+\ell)$ bits that returns a correct answer or $\bot$, with $\Pr[P(M_1,\ldots,M_k)=\bot]\le1/100$ for every input.
\end{authorclaim}
\begin{proof}
$P$ executes a randomized protocol for $\oplus_k\mathrm{BPHP}_{2^\ell}^N$ and verifies its answer by transferring additional $O(\ell)$ bits. The probability of failure can be reduced by repetition.
\end{proof}
We describe a protocol for $\mathrm{UDISJ}_{k,2^{\ell-k}-1}$ using $P$ for $\mathrm{ROW}\oplus_k\mathrm{BPHP}_{2^\ell}^{2^\ell+2^k}$. Inputs are $x_1,\ldots,x_k\in\{0,1\}^{2^{\ell-k}-1}$; $x_{i,j}$ is bit $j$ of $x_i$ and $\vec x=(x_1,\ldots,x_k)$.
\paragraph{Important matrices}
Let $\gamma:[2^{\ell-k}-1]\cup\{*\}\to\{0,1\}^{\ell-k}$ be bijective. Define matrices $D_1(x_1,\gamma),D_2(x_2),\ldots,D_k(x_k)$ of size $(2^\ell+2^k)\times\ell$. Let $A_i(b)$ be from Lemma21; for every $t\in\{0,1\}^{\ell-k}$ let $J_t$ have size $2^k\times(\ell-k)$ with every row $t$. Fix $W\in S_k$. Define
\[
D_1(x_1,\gamma):=\begin{pmatrix}
J_{\gamma(1)}&A_1(x_{1,1})\\\vdots&\vdots\\J_{\gamma(j)}&A_1(x_{1,j})\\\vdots&\vdots\\J_{\gamma(2^{\ell-k}-1)}&A_1(x_{1,2^{\ell-k}-1})\\J_{\gamma(*)}&W\\J_{\gamma(*)}&W
\end{pmatrix};
\]
for $i\in[k]\setminus\{1\}$,
\[
D_i(x_i):=\begin{pmatrix}
0_{2^k\times(\ell-k)}&A_i(x_{i,1})\\\vdots&\vdots\\0_{2^k\times(\ell-k)}&A_i(x_{i,j})\\\vdots&\vdots\\0_{2^k\times(\ell-k)}&A_i(x_{i,2^{\ell-k}-1})\\0_{2^k\times(\ell-k)}&0_{2^k\times k}\\0_{2^k\times(\ell-k)}&0_{2^k\times k}
\end{pmatrix}.
\]
The last $2^{k+1}$ rows of $D_1(x_1,\gamma)$ contain each row exactly twice.
Define $H_{\vec x}(\gamma):=D_1(x_1,\gamma)+D_2(x_2)+\cdots+D(x_k)$. Lemma21 gives the following property in the standard basis.
\begin{authorclaim}{Definition 23}
Let $M\in\mathbb F_2^{(2^k+2^\ell)\times\ell}$, $\gamma:[2^{\ell-k}-1]\cup\{*\}\to\{0,1\}^{\ell-k}$ bijective, and $e_1,\ldots,e_\ell$ a basis in $\mathbb F^\ell$. $M$ satisfies the key property w.r.t. $(\gamma,\vec x)$ in this basis if:
\begin{itemize}
\item For a row $s$ among the last $2^{k+1}$ rows, the first $\ell-k$ basis coordinates are $\gamma(*)_1,\ldots,\gamma(*)_{\ell-k}$ and $s$ appears exactly twice.
\item For a row $s$ with index in $[2^k(i-1)+1;2^ki]$, $i\in[2^{\ell-k}-1]$, the first $\ell-k$ basis coordinates are $\gamma(i)_1,\ldots,\gamma(i)_{\ell-k}$. If $\bigwedge_{j=1}^kx_{i,j}=0$ it appears exactly once; if $\bigwedge_{j=1}^kx_{i,j}=1$ it appears exactly twice and its $(\ell-k+1)$th basis coordinate is 0.
\end{itemize}
\end{authorclaim}
Consider an invertible matrix $E\in\mathbb F_2^{\ell\times\ell}$. Let $e_1,e_2,\ldots,e_\ell$ be the rows of $E$. Since $E$ is invertible, $e_1,e_2,\ldots,e_\ell$ form a basis. Let us define $C_{\vec x}(\gamma,E):=H(\vec x,\gamma)E$. Rows of $C_{\vec x}(\gamma,E)$ can be viewed as vectors with coordinates in the basis $e_1,e_2,\ldots,e_\ell$ corresponding to the rows of $H(\vec x,\gamma)$. Hence, $C_{\vec x}(\gamma,E)$ satisfies the key property w.r.t. $(\gamma,\vec x)$ in the basis $(e_1,e_2,\ldots,e_\ell)$.
Define $\operatorname{Fake}(\gamma,E)\subseteq\mathbb F_2^\ell$ to be the set of the last $2^{k+1}$ rows of $C_{\vec x}(\gamma,E)$. This set does not depend on $\vec x$. Its rows appear exactly twice.
\paragraph{Random variables}
Random variables are bold. Public randomness consists of:
\begin{itemize}
\item Uniform random bijection $\boldsymbol\gamma:[2^{\ell-k}-1]\cup\{*\}\to\{0,1\}^{\ell-k}$.
\item Uniform random invertible matrix $\mathbf E\in\mathbb F_2^{\ell\times\ell}$.
\item Uniform random permutation $\boldsymbol\pi$ of $[2^\ell+2^k]$, with corresponding permutation matrix $M_{\boldsymbol\pi}$: $(M_{\boldsymbol\pi})_{i,j}=1\Longleftrightarrow\boldsymbol\pi(i)=j$.
\item Uniform random matrices $\boldsymbol\Delta_1,\ldots,\boldsymbol\Delta_k\in\mathbb F_2^{(2^\ell+2^k)\times\ell}$ subject to $\sum_i\boldsymbol\Delta_i=0$.
\end{itemize}
Define $\mathbf P_i=M_{\boldsymbol\pi}D_i(x_i)\mathbf E+\boldsymbol\Delta_i$ for $i\ge2$ and $\mathbf P_1=M_{\boldsymbol\pi}D_1(x_1,\boldsymbol\gamma)\mathbf E+\boldsymbol\Delta_1$. The addition of $\boldsymbol\Delta_i$ makes each $\mathbf P_i$ indistinguishable from a random matrix. Moreover, $\sum_i\mathbf P_i=M_{\boldsymbol\pi}C_{\vec x}(\boldsymbol\gamma,\mathbf E)$, obtained by permuting its rows. Let $N$ be a constant to be chosen later.
\begin{authorclaim}{Algorithm 1: Protocol T solving $\mathrm{UDISJ}_{k,2^{\ell-k}-1}$}
Input $x_1,\ldots,x_k\in\{0,1\}^{2^{\ell-k}-1}$; $x_i$ is on party $i$'s forehead.
\begin{algorithmic}
\State $Z:=\varnothing$
\Loop\Statex repeat $N$ times
\State Sample $\pi\leftarrow\boldsymbol\pi$, $E\leftarrow\mathbf E$, $\gamma\leftarrow\boldsymbol\gamma$, $\vec\Delta\leftarrow\boldsymbol{\vec\Delta}$ \Comment{Fresh public bits}
\State $P_1:=M_\pi D_1(x_1,\gamma)E+\Delta_1$
\Statex $\triangleright$ Computable by parties $2,\ldots,k$.
\State $P_i:=M_\pi D_i(x_i)E+\Delta_i$ for $i\ge2$
\Statex $\triangleright$ Computable by all parties except $i$.
\State $z:=P(P_1,\ldots,P_k)$
\Statex $\triangleright$ Fresh random bits for $P$; $P_i$ is on party $i$'s forehead.
\State $Z:=Z\cup\{z\}$
\EndLoop
\If{$Z=\{\bot\}$}\State \Return $\bot$
\ElsIf{$Z\setminus\{\bot\}\subseteq\operatorname{Fake}(\gamma,E)$}\State \Return 1
\Statex $\triangleright$ Most likely there are no repeated rows except Fake; the key property gives $\mathrm{DISJ}(x_1,\ldots,x_k)=1$.
\EndIf
\State \Return 0
\end{algorithmic}
\end{authorclaim}
\paragraph{Protocol analysis}
Let us analyze the protocol $T$. Since it executes the protocol $P$ a constant number of times, $T$ transmits $O(t+\ell)$ bits. Assume that $x_1,x_2,\ldots,x_k$ is a 1-instance of $\mathrm{UDISJ}_{k,2^\ell+2^k}$. Then by the key property of $C_{\vec x}(\boldsymbol\gamma,\mathbf E)$ all repeated rows of $\sum_{i=1}^k\mathbf P_i$ are in $\operatorname{Fake}(\boldsymbol\gamma,\mathbf E)$, hence the protocol $T$ returns either $\bot$ or the correct answer. Since $P$ is executed $N$ times independently, the probability that $Z=\{\bot\}$ is at most $1/100^N$, hence $T$ returns1 with probability at least $1-1/100^N$.
The rest of the proof is devoted to the analysis of the case, where $x_1,x_2,\ldots,x_k$ is a 0-instance of $\mathrm{UDISJ}_{k,2^\ell+2^k}$. This is the most technically involved part of the proof. So it is a good point to give a large scale overview of the further proof strategy. Our goal is to show that if $x_1,x_2,\ldots,x_k$ is a 0-instance of $\mathrm{UDISJ}_{k,2^\ell+2^k}$, then the probability that $P(\mathbf P_1,\ldots,\mathbf P_k)$ returns a value from $\operatorname{Fake}(\boldsymbol\gamma,\mathbf E)$ is bounded by some constant less than1.
The random variable $P(\mathbf P_1,\ldots,\mathbf P_k)$ depends on random bits used by the protocol $P$ and on random bits needed for sampling $\mathbf P_1,\ldots,\mathbf P_k$. Let $R$ denote the set of all random strings used by the protocol $P$ (i.e. we assume that $P$ sample a random string from $R$ and use it as public randomness) and $S$ denote the set of all random strings used for sampling $\mathbf P_1,\ldots,\mathbf P_k$. We would like to construct two bijections $\alpha$ and $\beta$ on the set $S$ such that for every $s\in S$ the following two properties hold.
\begin{enumerate}
\item The three values of random variable $(\mathbf P_1,\ldots,\mathbf P_k)$ sampled using three strings $s,\alpha(s)$ and $\beta(s)$ as a random source, are the same.
\item Let $(\gamma,E)$, $(\gamma_\alpha,E_\alpha)$ and $(\gamma_\beta,E_\beta)$ be values of the random variable $(\boldsymbol\gamma,\mathbf E)$ that is sampled using three strings $s,\alpha(s)$ and $\beta(s)$ as a random source. Then $\operatorname{Fake}(\gamma,E)\cap\operatorname{Fake}(\gamma_\alpha,E_\alpha)\cap\operatorname{Fake}(\gamma_\beta,E_\beta)=\varnothing$.
\end{enumerate}
Consider arbitrary strings $r\in R$ and $s\in S$. The first property implies that for random variables sampled using strings $(r,s)$, $(r,\alpha(s))$ and $(r,\beta(s))$ as a random source values of $P(\mathbf P_1,\ldots,\mathbf P_k)$ are the same. The second property implies that for at least one of this cases this value does not belong to $\operatorname{Fake}(\boldsymbol\gamma,\mathbf E)$. Then, using that $\alpha$ and $\beta$ are bijections, we get $\Pr[P(\mathbf P_1,\ldots,\mathbf P_k)\in\operatorname{Fake}(\boldsymbol\gamma,\mathbf E)]\le2/3$.
Since we have many random variables, it is a tedious task to construct such $\alpha$ and $\beta$. In order to simplify this task we slightly relax the properties. We will define bijections $\alpha$ and $\beta$ not on all strings $S$ but only on the part of bits corresponding to sampling of $\boldsymbol\gamma$ and $\mathbf E$. More precisely we will define two bijections $\alpha$ and $\beta$ on the set of values of the random variable $(\boldsymbol\gamma,\mathbf E)$. We relax the first property as follows:
\begin{enumerate}
\item[1'.] For every $\gamma$ and $E$ the three conditional distributions of the random variable $(\mathbf P_1,\ldots,\mathbf P_k)$ under the following three conditions coincide:
\begin{enumerate}\item $(\boldsymbol\gamma,\mathbf E)=(\gamma,E)$,\item $(\boldsymbol\gamma,\mathbf E)=\alpha(\gamma,E)$ and\item $(\boldsymbol\gamma,\mathbf E)=\beta(\gamma,E)$.\end{enumerate}
\end{enumerate}
Unfortunately, we were not able to construct such bijections on the set of all pairs $(\gamma,E)$. Thus we take a set $\Xi$ consisting $1-\delta$ fraction of all values of $(\boldsymbol\gamma,\mathbf E)$ and we will claim that $\alpha$ and $\beta$ are bijections on $\Xi$. Such relaxation will weaken the bound of the probability up to $2/3+\delta$. We formalize the requirements to $\Xi,\alpha$ and $\beta$ in Definition24. Then we verify in Claim25 that these requirements are sufficient to bound $\Pr[P(\mathbf P_1,\ldots,\mathbf P_k)\in\operatorname{Fake}(\boldsymbol\gamma,\mathbf E)]$. The construction of $\Xi,\alpha$ and $\beta$ is given in Subsection5.3.
\begin{authorclaim}{Definition 24}
For a 0-instance $\vec x$ of $\mathrm{UDISJ}_{k,2^{\ell-k}-1}$ and constant $1>\delta\ge0$, let $\Xi$ consist of pairs $(\gamma,E)$ with $\gamma:[2^{\ell-k}-1]\cup\{*\}\to\{0,1\}^{\ell-k}$ bijective and $E\in\mathbb F_2^{\ell\times\ell}$ invertible. Bijections $\alpha,\beta:\Xi\to\Xi$ form a $(1-\delta)$-symmetry randomness space if:
\begin{itemize}
\item (Largeness) $\Pr[(\boldsymbol\gamma,\mathbf E)\in\Xi]\ge1-\delta$.
\item (Difference) $\operatorname{Fake}(\gamma,E)\cap\operatorname{Fake}(\alpha(\gamma,E))\cap\operatorname{Fake}(\beta(\gamma,E))=\varnothing$ for every $(\gamma,E)\in\Xi$.
\item (Symmetry) $C_{\vec x}(\gamma,E)$, $C_{\vec x}(\alpha(\gamma,E))$, $C_{\vec x}(\beta(\gamma,E))$ differ only by row permutation.
\end{itemize}
\end{authorclaim}
\begin{authorclaim}{Claim 25}
Under Definition24, $\Pr[P(\mathbf P_1,\ldots,\mathbf P_k)\in\operatorname{Fake}(\boldsymbol\gamma,\mathbf E)]\le2/3+\delta$.
\end{authorclaim}
\begin{proof}
Write $\boldsymbol{\vec P}=(\mathbf P_1,\ldots,\mathbf P_k)$, $\boldsymbol{\vec\Delta}=(\boldsymbol\Delta_1,\ldots,\boldsymbol\Delta_k)$, $\vec D(\vec x,\gamma)=(D_1(x_1,\gamma),D_2(x_2),\ldots,D_k(x_k))$, so $\boldsymbol{\vec P}=\boldsymbol{\vec\Delta}+M_{\boldsymbol\pi}(\vec D(\vec x,\boldsymbol\gamma)\mathbf E)$.
Let $p=\Pr[P(\boldsymbol{\vec P})\in\operatorname{Fake}(\boldsymbol\gamma,\mathbf E)]$. Then
\begin{align*}
p&=\sum_{\gamma,E}\Pr[P(\boldsymbol{\vec\Delta}+M_{\boldsymbol\pi}(\vec D(\vec x,\gamma)E))\in\operatorname{Fake}(\gamma,E)]\Pr[\boldsymbol\gamma=\gamma,\mathbf E=E]\\
&\le\sum_{(\gamma,E)\in\Xi}\Pr[P(\boldsymbol{\vec\Delta}+M_{\boldsymbol\pi}(\vec D(\vec x,\gamma)E))\in\operatorname{Fake}(\gamma,E)]\Pr[\boldsymbol\gamma=\gamma,\mathbf E=E]+\delta .
\end{align*}
Notice that for fixed $\gamma,E$ the random variable $\boldsymbol{\vec\Delta}+M_{\boldsymbol\pi}\cdot(\vec D(\vec x,\gamma)\cdot E)$ is distributed uniformly on the set of tuples $(L_1,\ldots,L_k)$ of $k$ matrices from $\mathbb F_2^{(2^\ell+2^k)\times\ell}$ such that $\sum_{i=1}^kL_i$ differs from $C_{\vec x}(\gamma,E)$ only by a permutation of rows. Let $(\gamma_{\alpha^{-1}},E_{\alpha^{-1}})=\alpha^{-1}(\gamma,E)$. By the symmetry condition, matrices $C_{\vec x}(\gamma,E)$ and $C_{\vec x}(\gamma_{\alpha^{-1}},E_{\alpha^{-1}})$ differ only by permutation of rows. Thus, for every set $A$ the probability
\[
\Pr[P(\boldsymbol{\vec\Delta}+M_{\boldsymbol\pi}(\vec D(\vec x,\gamma)E))\in A]
=\Pr[P(\boldsymbol{\vec\Delta}+M_{\boldsymbol\pi}(\vec D(\vec x,\gamma_{\alpha^{-1}})E_{\alpha^{-1}}))\in A].
\]
Therefore
\begin{align*}
p&\le\sum_{(\gamma,E)\in\Xi}\Pr[P(\boldsymbol{\vec\Delta}+M_{\boldsymbol\pi}(\vec D(\vec x,\gamma_{\alpha^{-1}})E_{\alpha^{-1}}))\in\operatorname{Fake}(\gamma,E)]\Pr[\boldsymbol\gamma=\gamma,\mathbf E=E]+\delta\\
&=\sum_{(\gamma,E)\in\Xi}\Pr[P(\boldsymbol{\vec\Delta}+M_{\boldsymbol\pi}(\vec D(\vec x,\gamma)E))\in\operatorname{Fake}(\alpha(\gamma,E))]\Pr[(\boldsymbol\gamma,\mathbf E)=\alpha(\gamma,E)]+\delta\\
&=\sum_{(\gamma,E)\in\Xi}\Pr[P(\boldsymbol{\vec\Delta}+M_{\boldsymbol\pi}(\vec D(\vec x,\gamma)E))\in\operatorname{Fake}(\alpha(\gamma,E))]\Pr[(\boldsymbol\gamma,\mathbf E)=(\gamma,E)]+\delta\\
&=\Pr[P(\boldsymbol{\vec P})\in\operatorname{Fake}(\alpha(\boldsymbol\gamma,\mathbf E)),(\boldsymbol\gamma,\mathbf E)\in\Xi]+\delta .
\end{align*}
Analogously, $p\le\Pr[P(\boldsymbol{\vec P})\in\operatorname{Fake}(\beta(\boldsymbol\gamma,\mathbf E)),(\boldsymbol\gamma,\mathbf E)\in\Xi]+\delta$. Also the inequality $p\le\Pr[P(\boldsymbol{\vec P})\in\operatorname{Fake}(\boldsymbol\gamma,\mathbf E),(\boldsymbol\gamma,\mathbf E)\in\Xi]+\delta$ follows by the largeness condition. Then,
\begin{align*}
3(1-p)&\ge\Pr[P(\boldsymbol{\vec P})\notin\operatorname{Fake}(\beta(\boldsymbol\gamma,\mathbf E))\vee(\boldsymbol\gamma,\mathbf E)\notin\Xi]\\
&\quad+\Pr[P(\boldsymbol{\vec P})\notin\operatorname{Fake}(\alpha(\boldsymbol\gamma,\mathbf E))\vee(\boldsymbol\gamma,\mathbf E)\notin\Xi]\\
&\quad+\Pr[P(\boldsymbol{\vec P})\notin\operatorname{Fake}(\boldsymbol\gamma,\mathbf E)\vee(\boldsymbol\gamma,\mathbf E)\notin\Xi]-3\delta\\
&\ge\Pr[P(\boldsymbol{\vec P})\notin\operatorname{Fake}(\boldsymbol\gamma,\mathbf E)\cap\operatorname{Fake}(\alpha(\boldsymbol\gamma,\mathbf E))\cap\operatorname{Fake}(\beta(\boldsymbol\gamma,\mathbf E))\vee(\boldsymbol\gamma,\mathbf E)\notin\Xi]-3\delta\\
&=1-3\delta .
\end{align*}
The last equality follows by the difference condition. Hence, $3(1-p)\ge1-3\delta$, thus $p\le2/3+\delta$.
\end{proof}
\begin{authorclaim}{Lemma 26}
For every 0-instance $\vec x$ of $\mathrm{UDISJ}_{k,2^{\ell-k}-1}$, for some $\delta<1/3-1/100$ there is a $(1-\delta)$-symmetry randomness space.
\end{authorclaim}
Lemma26 (proved in Section5.3) and Claim25 imply for a constant $\varepsilon>0$ that
\[
\Pr[P(\mathbf P_1,\ldots,\mathbf P_k)\notin\operatorname{Fake}(\boldsymbol\gamma,\mathbf E)]\ge\varepsilon+1/100,
\quad
\Pr[P(\mathbf P_1,\ldots,\mathbf P_k)\notin\operatorname{Fake}(\boldsymbol\gamma,\mathbf E)\cup\{\bot\}]\ge\varepsilon .
\]
Then, for $N=O(\log(1/\varepsilon))$, $T$ gives a correct answer for every 0-instance with probability at least $2/3$.
\end{proof}
